\textbf{Temporal Action Localization.}
This task (TAL) aims at localizing the start and end points of action clips from the entire untrimmed video with full observation.
We evaluate our InternVideo on four classic TAL datasets: THUMOS-14 \cite{idrees2017thumos}, ActivityNet-v1.3 \cite{caba2015activitynet}, HACS Segment  \cite{zhao2019hacs} and FineAction \cite{liu2022fineaction}.
In line with the previous temporal action localization tasks, we use mean Average Precision (mAP) for quantitative evaluations. Average Precision (AP) is calculated for each action category, to evaluate the proposals on action categories. It is computed under different tIoU thresholds.
We report the performance of the state-of-the-art TAL methods whose codes are publicly available, including ActionFormer \cite{zhang2022actionformer} method for THUMOS-14, ActivityNet-v1.3, and FineAction and TCANet \cite{qing2021temporal} method for HACS Segment.

We use ViT-H from our InternVideo as backbones for feature extraction. 
In our experiment, ViT-H models are pretrained from Hybrid datasets.
%As shown in Table \ref{tab:SOAT_th14}, Table \ref{tab:SOAT_anet},  Table \ref{tab:SOAT_hacs} and Table \ref{tab:SOAT_fineaction}.
As shown in Table \ref{tab:tal_emsemble},
our \model~outperforms best than all the preview methods on these four TAL datasets.
Note that, our InternVideo achieves huge improvements in temporal action localization, especially in fine-grained TAL datasets such as THUMOS-14 and FineAction.



